\documentclass{aip-cp}

\usepackage[numbers]{natbib}
\usepackage{rotating}
\usepackage{graphicx}

% Document starts
\begin{document}

% Title portion
\title{The Title Goes Here with Each Initial Letter Capitalized}

\author[aff1]{Konstantin Barkalov\corref{cor1}}
\author[aff1]{Ilya Lebedev}
\author[aff2]{Daniil Ryazanov}
\author[aff2]{Sergey Strijak}

\affil[aff1]{Lobachevsky State University of Nizhny Novgorod, Nizhny Novgorod, Russia}
\affil[aff2]{Ivannikov Institute for System Programming of the RAS, Moscow, Russia}
\corresp[cor1]{Corresponding author: konstantin.barkalov@itmm.unn.ru}

\maketitle

\begin{abstract}
The AIP Proceedings article template has many predefined paragraph styles for you to use/apply as you write your paper. To format your abstract, use the \LaTeX template style: {\itshape Abstract.} Each paper must include an abstract. Begin the abstract with the word ``Abstract'' followed by a period in bold font, and then continue with a normal 9 point font.
\end{abstract}

% Head 1
\section{INTRODUCTION}


\section{GLOBAL SEARCH ALGORITHM}

$y=(y_1,y_2,...,y_N)$

\begin{eqnarray}\label{problem}
&\varphi(y^\ast)=\min{\left\{\varphi(y):y\in D\right\}},\\
&D=\left\{y\in R^N: a_j\leq y_j \leq b_j, 1\leq j \leq N \right\}.\label{D}
\end{eqnarray}

$y^*$,

$y^i\in D, \; 1 \leq i \leq k$.

$\varphi(y)$

\[
\left|\varphi(y_1)-\varphi(y_2)\right|\leq L\left\|y_1-y_2\right\|,\; y_1,y_2 \in D, 0<L<\infty.
\]

\[
\varphi(y^\ast)=\varphi(y(x^\ast))=\min{\left\{\varphi(y(x)): x\in[0,1]\right\}},
\]

$x^1,...,x^k$

\begin{equation}\label{x_i}
0=x_0<x_1<\dots <x_k<x_{k+1}=1,
\end{equation}

$z_i=\varphi(y(x_i)), \; 0 \leq i\leq k+1,$ 

$(x_{i-1},x_i), \; 1\leq i \leq k, $


\[
R(t)=\max{R(i):1 \leq i \leq k}.
\]

$(x_{t-1},x_t)$, i.e. $x^{k+1} \in (x_{t-1},x_t)$.

  $\left(x_t-x_{t-1}\right)^{1/N}<\epsilon$
	

\section{NUMERICAL EXPERIMENTS}

\[
\nabla \cdot \vec{U} = 0
\]

\[
\frac{\partial \vec{U}}{\partial t} + \nabla \cdot \left( \vec{U} \otimes \vec{U}\right) - \nabla \cdot v \left( \nabla \vec{U} + (\nabla \vec{U})^T\right) = - \frac{1}{\rho} \nabla p.
\]

\[
\vec{U}(x,y) = 0, \; p(x,y) = 0.
\]

\[
\nabla p \cdot \vec{n} = 0, \; \vec{U} = 0.
\]

\[
\nabla p \cdot \vec{n} = 0, \; \vec{U} = \vec{U_i}.
\]

\[
 p = 0, \; \vec{U} = \vec{U_o}.
\]


\begin{figure}%[ht]
\includegraphics[width=1.0\linewidth]{fig1.png}
\caption{Figure 1}
\label{fig}
\end{figure}



Do not abbreviate Figure, Equation, etc.; display items are always singular, i.e., Figure 1 and 2. Equations are always singular, i.e., Equation 1 and 2, and should be inserted using the Equation Editor, not as graphics, in the main text.  Display items and captions should be inserted after the reference section. Please do not use footnotes in the text, additional information can be added to the reference list.

Cite all figures in the text consecutively. The word ``Figure'' should be spelled out if it is the first word of the sentence and abbreviated as ``Fig.'' elsewhere in the text. Place the figures as close as possible to their first mention in the text at the top or bottom of the page with the figure caption positioned below the figure, all centered. Figures must be inserted in the text and may not follow the Reference section. Set figure captions in 9 point size, Times Roman font. Type the word ``FIGURE 1.'' in bold uppercase, followed by a period.\footnote{This is an example of a footnote.}

% Acknowledgement
\section{ACKNOWLEDGMENTS}
This study was supported by the Russian Science Foundation, project No 21-11-00204.

\cite{Strongin2000,Kalyulin2017,Evtushenko2013,Jones2009,Zilinskas2010,Pinter1996,globalizerSystem,Barkalov2010,Strongin2020}



% References

%\nocite{*}
\bibliographystyle{aipnum-cp}%
\bibliography{bibliography}%


\end{document}
